\section{Conclusão}

\subsection{Síntese dos achados}

A tendência de alta de TDS identificada por ajuste linear global (2,9 a 3,9~mg/L/ano, p$<$0,01) não é robusta a recortes temporais nem a certas correções de autocorrelação -- a leitura mais honesta dos dados é que a série responde a um \textbf{padrão cíclico ligado a ciclos de seca da Califórnia}, não a uma tendência monotônica. Essa hipótese foi confirmada com uma variável climática externa e independente (PDSI, cujos \textit{changepoints}, detectados sem informação das datas do TDS, coincidem com as viradas de regime observadas) e decomposta em duas frentes causais: o WRTDS mostra que a tendência flow-normalized não é significativa (p=0,064) -- quase toda a tendência bruta é efeito de vazão; e o balanço de massa confirma diretamente que a carga de sal caiu ($-3{,}20\%$/ano) enquanto a vazão caiu mais rápido ainda ($-3{,}85\%$/ano) -- \textbf{diluição, não mais sal entrando no sistema}. Como a extrapolação linear desses componentes cruza valores fisicamente impossíveis em horizontes longos, a previsão de +20 anos é apresentada como uma faixa condicionada a cenários climáticos (594 a 718~mg/L), mais estreita e mais baixa que a convergência dos três finalistas de previsão pontual ($\sim$760 a 800~mg/L), porque limita a vazão futura a valores historicamente plausíveis.

TDS correlaciona-se positivamente com Amônia mesmo destendenciado ($r=0{,}175$, p=0,018), oferecendo evidência de que salinidade elevada coincide com nitrificação reduzida. A correlação com BOD é negativa e não significativa em todos os tratamentos testados -- o oposto do previsto, reportado sem suavização; a leitura mais provável é a falta de variação real na série de BOD, rente ao limite de detecção durante quase todo o período.

\subsection{Interpretação em contexto real}

Os achados são compatíveis com o mecanismo descrito por Schwabe et al.\ \cite{schwabe2020unintended} (conservação $\to$ menos vazão $\to$ mais concentrado) e com Wolfand et al.\ \cite{wolfand2022dilution}, que encontraram o mesmo padrão na bacia do rio Los Angeles. Mas um relatório técnico independente com 26 estações da região \cite{scsc2018tds}, cujo texto completo foi lido na íntegra, atribui a maior parte da variabilidade do TDS de influente (~88\%) ao TDS da própria água de abastecimento -- não à conservação local (~12\%) -- sugerindo que o mecanismo de Schwabe et al.\ é real, porém secundário. O balanço de massa deste trabalho corrobora essa leitura de forma independente e local: a carga de sal na LAGWRP não está subindo, está caindo, e a vazão cai mais rápido ainda.

\subsection{Limitações}

Este é um estudo de ponto único, não generalizável sem cautela; a análise de correlação é observacional, não estabelece causalidade; a série de BOD é fortemente censurada (65\% ND) e tem amplitude pequena, limitando o poder estatístico da correlação TDS--BOD; apenas 1 dos 35 métodos atinge $R^2$ positivo no \textit{holdout} de 24 meses, e o próprio \textit{baseline} naive supera a maioria dos métodos sofisticados nesse critério de curto prazo; a variável que a literatura da região aponta como mais explicativa -- o TDS da própria água de abastecimento (~88\% da variabilidade, \cite{scsc2018tds}) -- só foi parcialmente reconstruída (lado Los Angeles/LADWP, $r=0{,}912$, dado anual não mensal), não estando disponível como covariável direta na bateria de previsão original.

\subsection{Trabalhos futuros}

O escopo de dados deste trabalho está encerrado. A direção de maior potencial de impacto para trabalhos subsequentes é obter uma série mensal completa de TDS da água de abastecimento na origem (incluindo o lado Glendale, ainda sem dado equivalente), permitindo treinar os modelos com o principal determinante hipotético do TDS como covariável direta, em vez de inferi-lo apenas indiretamente via vazão reconstruída -- o que provavelmente reduziria de forma substancial a incerteza dos horizontes de longo prazo reportados aqui.
